% =============================================================================
% APPENDIX HEC: LIT-JAMES: James Heckman Research Integration
% =============================================================================
% Category: LIT
% Language: English
% Version: 1.0 (January 2026)
% Status: Draft
%
% This appendix integrates research papers by James Heckman into the EBF framework.
%
% =============================================================================

\chapter{LIT-JAMES: James Heckman Research Integration}
\label{app:hec}

\begin{abstract}
This appendix integrates the research contributions of James Heckman into the
Evidence-Based Framework. It documents key papers, core findings, and their
integration with EBF dimensions including complementarity parameters ($\gamma$),
context factors ($\Psi$), and utility dimensions.
\end{abstract}

\begin{tcolorbox}[colback=blue!5!white, colframe=blue!75!black,
    title=Quick Reference: James Heckman Research]

\textbf{Appendix:} HEC (LIT)

\textbf{Researcher:} James Heckman

\textbf{Core Themes:}
\begin{itemize}[nosep]
\item Behavioral economics foundations
\item Empirical evidence for EBF parameters
\item Policy implications and interventions
\end{itemize}

\textbf{EBF Integration:}
\begin{itemize}[nosep]
\item Complementarity values ($\gamma$) from experimental evidence
\item Context effects ($\Psi$) documented across studies
\item Utility dimension weightings calibrated to findings
\end{itemize}

\textbf{Cross-References:} BBB (CORE-WHERE), K (LIT-FEHR), G (Glossary)
\end{tcolorbox}

% -----------------------------------------------------------------------------
% APPENDIX SCOPE BOX
% -----------------------------------------------------------------------------

\begin{tcolorbox}[
    colback=orange!5!white,
    colframe=orange!75!black,
    title={\textbf{Appendix Scope: Ziel / In-Scope / Out-of-Scope / Constraints}},
    fonttitle=\bfseries
]

\textbf{Ziel (Objective):}
\begin{itemize}[nosep]
    \item Integrate James Heckman's research into EBF with parameter extraction
\end{itemize}

\textbf{In-Scope:}
\begin{itemize}[nosep]
    \item Paper-by-paper integration with 10C mapping
    \item Extraction of $\gamma$, $\Psi$, and effect size parameters
    \item Cross-references to related EBF components
\end{itemize}

\textbf{Out-of-Scope (delegiert an andere Appendices):}
\begin{itemize}[nosep]
    \item Parameter validation $\rightarrow$ Appendix UNMAPPED_BBB (CORE-WHERE)
    \item Formal axiomatization $\rightarrow$ CORE appendices
    \item Meta-analysis $\rightarrow$ LIT-M appendices
\end{itemize}

\textbf{Constraints (Anwendungsgrenzen):}
\begin{itemize}[nosep]
    \item Parameters are extracted from published studies
    \item Effect sizes may vary across replications
    \item Context-specificity of findings applies
\end{itemize}

\textbf{Lieferobjekte (Deliverables):}
\begin{itemize}[nosep]
    \item Paper integration table with 10C mappings
    \item Extracted parameters for BBB repository
    \item Cross-reference map to related literature
\end{itemize}

\end{tcolorbox}

%%%%%%%%%%%%%%%%%%%%%%%%%%%%%%%%%%%%%%%%%%%%%%%%%%%%%%%%%%%%%%%%%%%%%%%%%%%%%%%
\section{The Fundamental Question}
\label{sec:hec-fundamental}
%%%%%%%%%%%%%%%%%%%%%%%%%%%%%%%%%%%%%%%%%%%%%%%%%%%%%%%%%%%%%%%%%%%%%%%%%%%%%%%

\begin{tcolorbox}[
    colback=red!5!white,
    colframe=red!75!black,
    title={\textbf{Fundamental Question}}
]
\textbf{How does lit-james: james heckman research integration integrate with and contribute to the
Evidence-Based Framework for understanding behavioral outcomes?}

This appendix addresses the core challenge of formalizing behavioral principles
within a rigorous theoretical framework while maintaining practical applicability.
\end{tcolorbox}




\section{The Human Capital Foundation: Twenty Papers by James Heckman}
\label{app:heckman}
%%%%%%%%%%%%%%%%%%%%%%%%%%%%%%%%%%%%%%%%%%%%%%%%%%%%%%%%%%%%%%%%%%%%%%%%%%%%%%%

This appendix integrates twenty of the most influential papers by James Heckman
(2000 Nobel laureate, John Bates Clark Medal 1983), one of the most cited
economists in history. His work spans econometric methodology (selection bias,
treatment effects) and human development (skill formation, early childhood).
Together with Fehr (behavioral), Acemoglu (political institutions), and
Shleifer (legal/financial), Heckman provides the \emph{human capital} pillar
of the $(C, \Psi)$ framework.

\subsection{Overview: Heckman's Research in Framework Terms}

Heckman's central insight concerns the \emph{technology of skill formation}:
\begin{equation}
  \theta_{t+1} = f(\theta_t, I_t, \Psi_t)
\end{equation}
where skills at $t+1$ depend on current skills $\theta_t$, investments $I_t$,
and context $\Psi_t$. This technology exhibits:
\begin{itemize}
  \item \textbf{Self-productivity:} $\partial\theta_{t+1}/\partial\theta_t > 0$
  \item \textbf{Dynamic complementarity:} $\partial^2\theta_{t+1}/\partial\theta_t\partial I_t > 0$
  \item \textbf{Sensitive periods:} $\partial f/\partial I_t$ varies with age
\end{itemize}

\begin{center}
\renewcommand{\arraystretch}{1.2}
\small
\begin{tabular}{ll}
\hline
\textbf{Framework Element} & \textbf{Heckman's Contribution} \\
\hline
$C^{skill,investment} > 0$ & Dynamic complementarity in skill formation \\
Self-productivity & Skills beget skills ($\theta_t \to \theta_{t+1}$) \\
Multiple skills & Cognitive + noncognitive dimensions \\
Selection bias & Observed $\neq$ potential outcomes \\
Treatment effects & Heterogeneous responses to interventions \\
Early investment returns & High ROI for disadvantaged children \\
\hline
\end{tabular}
\end{center}

\subsection{Part I: Econometric Methodology (9 Papers)}

These papers established the modern framework for causal inference in
observational data, addressing selection bias and treatment effect heterogeneity.

%--- Paper 1 ---
\subsubsection{Sample Selection Bias as Specification Error (1979)}

\paragraph{Citation.}
Heckman, J.J. (1979). ``Sample Selection Bias as a Specification Error.''
\emph{Econometrica}, 47(1): 153-161.

\paragraph{Core Contribution.}
When observations are non-random (e.g., we observe wages only for workers),
OLS estimates are biased. The ``Heckman correction'' provides consistent
estimation via two-step procedure.

\paragraph{Framework Integration.}
\begin{itemize}
  \item \textbf{Selection as $\Psi$:} Who we observe depends on context
  \item \textbf{Potential vs. observed:} $Y^{observed} \neq Y^{potential}$
  \item \textbf{Methodological foundation:} Enables causal inference
\end{itemize}

\paragraph{Framework Significance.}
This is one of the most cited papers in economics (50,000+ citations).
It established that \emph{context determines selection}, and selection
affects what we can learn from data.

%--- Paper 2 ---
\subsubsection{Common Structure of Selection Models (1976)}

\paragraph{Citation.}
Heckman, J.J. (1976). ``The Common Structure of Statistical Models of
Truncation, Sample Selection and Limited Dependent Variables.''
\emph{Annals of Economic and Social Measurement}, 5(4): 475-492.

\paragraph{Core Contribution.}
Unified framework showing that truncation, selection, and limited
dependent variable models share common structure. Simple estimator proposed.

\paragraph{Framework Integration.}
\begin{itemize}
  \item \textbf{Unified structure:} Different problems, same mathematics
  \item \textbf{Latent variables:} Unobserved heterogeneity pervasive
  \item \textbf{Estimation strategy:} Control functions for selection
\end{itemize}

%--- Paper 3 ---
\subsubsection{Shadow Prices and Labor Supply (1974)}

\paragraph{Citation.}
Heckman, J.J. (1974). ``Shadow Prices, Market Wages, and Labor Supply.''
\emph{Econometrica}, 42(4): 679-694.

\paragraph{Core Contribution.}
Reservation wages (``shadow prices'') determine labor force participation.
Corner solutions in labor supply require modeling participation decision
separately from hours decision.

\paragraph{Framework Integration.}
\begin{itemize}
  \item \textbf{Two-stage decision:} Participate, then choose hours
  \item \textbf{Corner solutions:} $L = 0$ for many individuals
  \item \textbf{Heterogeneity:} Reservation wages vary across people
\end{itemize}

%--- Paper 4 ---
\subsubsection{Duration Analysis (1984)}

\paragraph{Citation.}
Heckman, J.J. \& Singer, B. (1984). ``A Method for Minimizing the Impact
of Distributional Assumptions in Econometric Models for Duration Data.''
\emph{Econometrica}, 52(2): 271-320.

\paragraph{Core Contribution.}
Nonparametric approach to unobserved heterogeneity in duration models.
Avoids arbitrary distributional assumptions that can bias estimates.

\paragraph{Framework Integration.}
\begin{itemize}
  \item \textbf{Duration as $C^*$:} Time in state depends on context
  \item \textbf{Unobserved heterogeneity:} $\Psi_{individual}$ varies
  \item \textbf{Nonparametric estimation:} Let data reveal distribution
\end{itemize}

%--- Paper 5 ---
\subsubsection{Structural Equations and Treatment Effects (2005)}

\paragraph{Citation.}
Heckman, J.J. \& Vytlacil, E.J. (2005). ``Structural Equations, Treatment
Effects, and Econometric Policy Evaluation.'' \emph{Econometrica}, 73(3): 669-738.

\paragraph{Core Contribution.}
Unified framework connecting structural econometrics with treatment effect
literature. Marginal treatment effect (MTE) as fundamental building block.

\paragraph{Framework Integration.}
\begin{itemize}
  \item \textbf{MTE:} Treatment effect for person indifferent to treatment
  \item \textbf{Policy-relevant parameters:} Different from ATE, ATT
  \item \textbf{Selection on gains:} People choose based on their returns
\end{itemize}

%--- Paper 6 ---
\subsubsection{Matching as Econometric Estimator (1997)}

\paragraph{Citation.}
Heckman, J.J., Ichimura, H., \& Todd, P.E. (1997). ``Matching as an
Econometric Evaluation Estimator.'' \emph{Review of Economic Studies}, 64(4): 605-654.

\paragraph{Core Contribution.}
Conditions under which matching estimators identify causal effects.
Bias from failure of conditional independence assumption.

\paragraph{Framework Integration.}
\begin{itemize}
  \item \textbf{Conditional independence:} Selection on observables
  \item \textbf{Common support:} Need comparable treated/untreated
  \item \textbf{Practical guidance:} When matching works, when it fails
\end{itemize}

%--- Paper 7 ---
\subsubsection{Econometric Causality: The Central Role of Thought Experiments (2024)}

\paragraph{Citation.}
Heckman, J.J. \& Pinto, R. (2024). ``Econometric Causality: The Central Role of
Thought Experiments.'' \emph{Journal of Econometrics}, 243(1): 105719.
\citep{heckman2024causality}

\paragraph{Core Contribution.}
Establishes thought experiments as the foundation of econometric causality,
contrasting the Frisch-Haavelmo tradition with Neyman-Rubin potential outcomes
and Pearl's Do-Calculus. Only structural econometric models can answer all
four fundamental policy questions (P1-P4).

\paragraph{Framework Integration.}
\begin{itemize}
  \item \textbf{Fixing vs. Conditioning:} Fixing modifies mechanisms; conditioning restricts samples
  \item \textbf{P1-P4 Framework:} Treatment on Treated (P1), Treatment on Untreated (P2),
        Policy Forecast (P3), New Policy Effects (P4)
  \item \textbf{Three Tasks:} Defining hypotheticals (T1), Identification (T2), Estimation (T3)
  \item \textbf{Structural Parameters:} $\lambda$, $\beta$, $\gamma$ are ``deep'' and policy-invariant
\end{itemize}

\paragraph{Framework Significance.}
\textbf{★★★★★ FOUNDATIONAL PAPER FOR EBF METHODOLOGY.}
This paper provides the epistemological justification for the entire EBF architecture:
\begin{enumerate}
  \item \textbf{Why structural parameters:} $\lambda$ (loss aversion), $\beta$ (discount factor),
        $\gamma$ (complementarity) are \emph{not} treatment effects but deep parameters
  \item \textbf{Why context ($\Psi$) matters:} Context = which mechanisms are held fixed
        in the thought experiment
  \item \textbf{Why P1-P4 framework:} Enables systematic intervention evaluation
  \item \textbf{Why 10C framework:} Each CORE dimension specifies a mechanism to fix or vary
\end{enumerate}

See also: Full-text archive at \texttt{data/paper-texts/PAP-heckman2024causality.md}

%--- Paper 8 ---
\subsubsection{Measuring the Growth of Skills (2026)}

\paragraph{Citation.}
Heckman, J.J., Tian, H., Zhang, Z. \& Zhou, J. (2026). ``Measuring the Growth of Skills.''
\emph{NBER Working Paper No. 34737}.
\citep{heckman2026measuring}

\paragraph{Core Contribution.}
Establishes that skills do not evolve as higher levels of a fixed set of abilities---new
skills \emph{emerge} with development, qualitatively different from earlier skills. Rejects
the common scale assumption for language and cognitive skills. Provides scale-free test
of dynamic complementarity without relying on psychometric fiat.

\paragraph{Framework Integration.}
\begin{itemize}
  \item \textbf{Skill Emergence:} New skills emerge at each developmental level---not just ``more of the same''
  \item \textbf{Common Scale Rejected:} Language and cognitive skills cannot be compared on common scale; fine motor CAN
  \item \textbf{Scale-Free Testing:} Dynamic complementarity testable using time-to-mastery within levels
  \item \textbf{Fadeout Explained:} Apparent fadeout may be measurement artifact of skill emergence
  \item \textbf{Catch-Up Effect:} $\gamma_{low} > \gamma_{normal} > \gamma_{high}$---disadvantaged benefit most
\end{itemize}

\paragraph{Key Parameters.}
\begin{itemize}
  \item \textbf{Dynamic Complementarity (Language/Cognitive):} Positive, significant at UHP levels 8, 9, 11
  \item \textbf{Dynamic Complementarity (Fine Motor):} NULL---no evidence
  \item \textbf{Time to First Mastery:} Full exposure children require significantly fewer trials
  \item \textbf{Common Scale Validity:} REJECTED for language/cognitive, NOT REJECTED for fine motor
\end{itemize}

\paragraph{Framework Significance.}
\textbf{★★★★ LEVEL 4 THEORY PAPER FOR EBF MEASUREMENT METHODOLOGY.}
\begin{enumerate}
  \item \textbf{Why scale matters:} EBF parameters ($\lambda$, $\beta$, $\gamma$) should not blindly
        assume cardinal psychometric scales across developmental stages
  \item \textbf{Why fadeout:} Apparent fadeout of early interventions may reflect skill emergence,
        not actual loss of treatment effects
  \item \textbf{Why within-level measurement:} Valid growth assessment requires measurement within
        developmental levels using level-specific scales
  \item \textbf{Why early investment:} Dynamic complementarity enables low-ability children to catch up
\end{enumerate}

Data: China REACH RCT with weekly assessments across 11 UHP difficulty levels.

See also: Full-text archive at \texttt{data/paper-texts/PAP-heckman2026measuring.md}

%--- Paper 9 ---
\subsubsection{The Economic Approach to Personality, Character and Virtue (2023)}

\paragraph{Citation.}
Heckman, J.J., Galaty, T. \& Tian, H. (2023). ``The Economic Approach to Personality,
Character and Virtue.'' \emph{NBER Working Paper No. 31258}.
\citep{heckman2023virtue}

\paragraph{Core Contribution.}
Develops unified economic framework for personality, character and virtue, bridging
economics, psychology, and philosophy. Key insight: \textbf{Traits are actions (behaviors),
not immutable genetic endowments}---personality is malleable through investment and own
effort. Formalizes Aristotelian virtue ethics (golden mean, habituation) in economic terms.

\paragraph{Framework Integration.}
\begin{itemize}
  \item \textbf{Traits as Actions:} $\theta_{t+1} = g(\theta_t, e_t, h_t)$---traits evolve
        through effort and situations
  \item \textbf{Character:} Stable dispositions producing consistent behavior across
        varied situations $h \in H$
  \item \textbf{Virtue:} Commitment to ideal point $\sigma^*$ (golden mean between extremes)
  \item \textbf{Self-Control Model:} $S_t$ as depletable resource for resisting temptation
  \item \textbf{Habituation Model:} $K_{t+1} = (1-\delta_K)K_t + h(a_t)$---practice builds
        character
  \item \textbf{Utility Structure:} $U(a, \psi, \theta, h)$ distinguishes preferences
        $\psi$ from traits $\theta$
\end{itemize}

\paragraph{Key Parameters (Theory MS-VE-001/002/003).}
\begin{itemize}
  \item \textbf{Self-Control Depreciation:} $\delta_S \in [0.10, 0.20]$ per period (ego
        depletion effect)
  \item \textbf{Habituation Depreciation:} $\delta_K \in [0.02, 0.05]$ per period (``use
        it or lose it'')
  \item \textbf{Cost Reduction:} $\text{Cost}(\text{virtue}) = c(a)/(1 + K_t)$---higher
        habituation $\to$ easier virtue
  \item \textbf{Golden Mean:} $\sigma^* = \arg\max U(\sigma)$ where $U(\theta) = -|\theta - \sigma^*|$
\end{itemize}

\paragraph{Framework Significance.}
\textbf{★★★★ LEVEL 4 FOUNDATIONAL BRIDGE PAPER.}
\begin{enumerate}
  \item \textbf{Economics $\leftrightarrow$ Psychology:} Unified framework for personality
        measurement, distinguishing traits from preferences operationally
  \item \textbf{Economics $\leftrightarrow$ Philosophy:} Formalizes Aristotle's virtue
        ethics (golden mean, habituation, akrasia as self-control failure)
  \item \textbf{Character Formation:} Converting effortful virtue into habitual virtue
        through practice
  \item \textbf{Intervention Design:} Perry/Jamaica/ABC/China REACH evidence shows
        character skills are malleable with high returns
\end{enumerate}

Aristotle: ``We become just by doing just acts.'' Economic formalization: $K_{t+1} = (1-\delta_K)K_t + h(a_t)$

See also: Full-text archive at \texttt{data/paper-texts/PAP-heckman2023virtue.md}

\subsection{Part II: Skill Formation and Human Development (7 Papers)}

%--- Paper 7 ---
\subsubsection{The Technology of Skill Formation (2007)}

\paragraph{Citation.}
Cunha, F. \& Heckman, J.J. (2007). ``The Technology of Skill Formation.''
\emph{American Economic Review}, 97(2): 31-47.

\paragraph{Core Contribution.}
Formal model of skill formation with self-productivity and dynamic
complementarity. Skills at $t$ raise productivity of investment at $t+1$.

\paragraph{Framework Integration.}
\begin{itemize}
  \item \textbf{Dynamic complementarity:} $C^{skill_t, investment_t} > 0$
  \item \textbf{Self-productivity:} Skills beget skills
  \item \textbf{Sensitive periods:} Investment timing matters
  \item \textbf{Multiple skills:} Cognitive and noncognitive
\end{itemize}

\paragraph{Framework Significance.}
This paper provides the microeconomic foundation for understanding why
early childhood investments have such high returns: dynamic complementarity
means early skills multiply the returns to later investments.

%--- Paper 8 ---
\subsubsection{Estimating Skill Formation Technology (2010)}

\paragraph{Citation.}
Cunha, F., Heckman, J.J., \& Schennach, S.M. (2010). ``Estimating the
Technology of Cognitive and Noncognitive Skill Formation.''
\emph{Econometrica}, 78(3): 883-931.

\paragraph{Core Finding.}
Structural estimation of skill formation technology. Self-productivity
and dynamic complementarity empirically confirmed. Noncognitive skills
more malleable than cognitive skills at later ages.

\paragraph{Framework Integration.}
\begin{itemize}
  \item \textbf{Empirical validation:} $C^{skill,investment} > 0$ confirmed
  \item \textbf{Differential malleability:} Noncognitive more plastic later
  \item \textbf{Policy implication:} Different interventions for different ages
\end{itemize}

%--- Paper 9 ---
\subsubsection{Schools, Skills and Synapses (2008)}

\paragraph{Citation.}
Heckman, J.J. (2008). ``Schools, Skills and Synapses.''
\emph{Economic Inquiry}, 46(3): 289-324.

\paragraph{Core Contribution.}
Integration of economics with neuroscience. Early experiences shape
brain architecture; skill gaps emerge before schooling begins.

\paragraph{Framework Integration.}
\begin{itemize}
  \item \textbf{$\Psi_{biology}$:} Neural development as context
  \item \textbf{Critical periods:} Brain plasticity declines with age
  \item \textbf{Pre-school gaps:} Inequality established before age 5
\end{itemize}

%--- Paper 10 ---
\subsubsection{Cognitive and Noncognitive Skills (2006)}

\paragraph{Citation.}
Heckman, J.J., Stixrud, J., \& Urzua, S. (2006). ``The Effects of Cognitive
and Noncognitive Abilities on Labor Market Outcomes and Social Behavior.''
\emph{Journal of Labor Economics}, 24(3): 411-482.

\paragraph{Core Finding.}
Both cognitive and noncognitive skills predict adult outcomes.
Noncognitive skills (conscientiousness, self-control) matter as much
as IQ for wages, employment, and social behavior.

\paragraph{Framework Integration.}
\begin{itemize}
  \item \textbf{Multidimensional human capital:} IQ + personality
  \item \textbf{$C^{cognitive,noncognitive}$:} Skills interact
  \item \textbf{Beyond IQ:} Character matters for success
\end{itemize}

%--- Paper 11 ---
\subsubsection{GED and Noncognitive Skills (2001)}

\paragraph{Citation.}
Heckman, J.J. \& Rubinstein, Y. (2001). ``The Importance of Noncognitive
Skills: Lessons from the GED Testing Program.''
\emph{American Economic Review}, 91(2): 145-149.

\paragraph{Core Finding.}
GED recipients have cognitive skills similar to high school graduates
but earn less. The difference: noncognitive skills. GED ``certifies''
cognitive ability but signals low noncognitive skills.

\paragraph{Framework Integration.}
\begin{itemize}
  \item \textbf{Noncognitive matters:} Persistence, self-control
  \item \textbf{Signaling:} GED reveals inability to complete HS
  \item \textbf{Policy failure:} Testing cognitive skills insufficient
\end{itemize}

%--- Paper 12 ---
\subsubsection{Perry Preschool Analysis (2010)}

\paragraph{Citation.}
Heckman, J.J., Moon, S.H., Pinto, R., Savelyev, P.A., \& Yavitz, A. (2010).
``The Rate of Return to the HighScope Perry Preschool Program.''
\emph{Journal of Public Economics}, 94(1-2): 114-128.

\paragraph{Core Finding.}
Perry Preschool program returned 7-10\% annually through age 40.
Returns came primarily through effects on noncognitive skills,
not IQ gains (which faded).

\paragraph{Framework Integration.}
\begin{itemize}
  \item \textbf{High ROI:} Early investment pays off
  \item \textbf{Mechanism:} Noncognitive skills, not IQ
  \item \textbf{Long-run effects:} Benefits persist to age 40+
\end{itemize}

%--- Paper 13 ---
\subsubsection{Early Childhood Investment (2013)}

\paragraph{Citation.}
Heckman, J.J. (2013). ``Giving Kids a Fair Chance.''
MIT Press (and related papers).

\paragraph{Core Contribution.}
Policy synthesis: invest early in disadvantaged children.
Cost-effective, produces equality and efficiency simultaneously.

\paragraph{Framework Integration.}
\begin{itemize}
  \item \textbf{No equity-efficiency tradeoff:} Early investment achieves both
  \item \textbf{Policy design:} Target disadvantaged families
  \item \textbf{Life-cycle perspective:} Early $\to$ middle $\to$ late
\end{itemize}

\subsection{Part III: Labor Markets and Policy Evaluation (7 Papers)}

%--- Paper 14 ---
\subsubsection{Economics of Active Labor Market Programs (1999)}

\paragraph{Citation.}
Heckman, J.J., LaLonde, R.J., \& Smith, J.A. (1999).
``The Economics and Econometrics of Active Labor Market Programs.''
\emph{Handbook of Labor Economics}, Vol. 3A, Chapter 31.

\paragraph{Core Finding.}
Comprehensive review finding most job training programs have modest effects.
Evaluation methodology matters enormously for conclusions.

\paragraph{Framework Integration.}
\begin{itemize}
  \item \textbf{Policy skepticism:} Many programs don't work
  \item \textbf{Heterogeneous effects:} Some subgroups benefit
  \item \textbf{Evaluation matters:} Method affects conclusions
\end{itemize}

%--- Paper 15 ---
\subsubsection{Impact of Federal Anti-Discrimination Policy (1989)}

\paragraph{Citation.}
Heckman, J.J. \& Payner, B.S. (1989). ``Determining the Impact of Federal
Antidiscrimination Policy on the Economic Status of Blacks.''
\emph{American Economic Review}, 79(1): 138-177.

\paragraph{Core Finding.}
Civil Rights Act of 1964 and federal enforcement improved black
economic status in South Carolina. Policy worked through changing
employer behavior.

\paragraph{Framework Integration.}
\begin{itemize}
  \item \textbf{$\Psi_{legal}$ change:} Anti-discrimination law
  \item \textbf{Enforcement matters:} De jure + de facto
  \item \textbf{Causal identification:} Timing, regional variation
\end{itemize}

%--- Paper 16 ---
\subsubsection{General Equilibrium Treatment Effects (1998)}

\paragraph{Citation.}
Heckman, J.J., Lochner, L., \& Taber, C. (1998). ``General Equilibrium
Treatment Effects: A Study of Tuition Policy.''
\emph{American Economic Review}, 88(2): 381-386.

\paragraph{Core Contribution.}
Partial equilibrium treatment effects miss general equilibrium responses.
Tuition subsidies change skill supplies, which changes wages.

\paragraph{Framework Integration.}
\begin{itemize}
  \item \textbf{General equilibrium:} Interventions change $\Psi$
  \item \textbf{Price effects:} More graduates $\to$ lower skill premium
  \item \textbf{Policy evaluation:} Partial effects misleading
\end{itemize}

%--- Paper 17 ---
\subsubsection{Rising Wage Inequality (1998)}

\paragraph{Citation.}
Heckman, J.J., Lochner, L., \& Taber, C. (1998). ``Explaining Rising
Wage Inequality: Explorations with a Dynamic General Equilibrium Model.''
\emph{Review of Economic Dynamics}, 1(1): 1-58.

\paragraph{Core Finding.}
Skill-biased technical change and cohort effects explain rising
wage inequality. Human capital investments respond to changing returns.

\paragraph{Framework Integration.}
\begin{itemize}
  \item \textbf{$\Psi_{tech}$ change:} Skill-biased technical change
  \item \textbf{Endogenous skills:} Investment responds to returns
  \item \textbf{Inequality dynamics:} Supply and demand for skills
\end{itemize}

%--- Paper 18 ---
\subsubsection{Empirical Content of the Roy Model (1990)}

\paragraph{Citation.}
Heckman, J.J. \& Honor\'e, B.E. (1990). ``The Empirical Content of the
Roy Model.'' \emph{Econometrica}, 58(5): 1121-1149.

\paragraph{Core Contribution.}
Roy model of self-selection: people choose sectors based on
comparative advantage. Identification without distributional assumptions.

\paragraph{Framework Integration.}
\begin{itemize}
  \item \textbf{Comparative advantage:} People sort by $\theta_{sector}$
  \item \textbf{Selection on returns:} Choose highest-return sector
  \item \textbf{Identification:} What can we learn from data?
\end{itemize}

%--- Paper 19 ---
\subsubsection{Family Labor Supply (1974)}

\paragraph{Citation.}
Ashenfelter, O. \& Heckman, J.J. (1974). ``The Estimation of Income and
Substitution Effects in a Model of Family Labor Supply.''
\emph{Econometrica}, 42(1): 73-85.

\paragraph{Core Contribution.}
Joint estimation of husband and wife labor supply. Family as
decision unit, not individual.

\paragraph{Framework Integration.}
\begin{itemize}
  \item \textbf{Family as unit:} $C^{husband,wife}$ in labor supply
  \item \textbf{Cross-effects:} Wife's wage affects husband's hours
  \item \textbf{Joint decisions:} Household optimization
\end{itemize}

%--- Paper 20 ---
\subsubsection{Economics and Psychology of Personality (2008)}

\paragraph{Citation.}
Borghans, L., Duckworth, A.L., Heckman, J.J., \& ter Weel, B. (2008).
``The Economics and Psychology of Personality Traits.''
\emph{Journal of Human Resources}, 43(4): 972-1059.

\paragraph{Core Contribution.}
Integration of personality psychology with economics. Big Five traits
predict economic outcomes. Measurement challenges and opportunities.

\paragraph{Framework Integration.}
\begin{itemize}
  \item \textbf{Personality as $\theta$:} Stable individual differences
  \item \textbf{Big Five:} Conscientiousness, openness, etc.
  \item \textbf{Predictive power:} Personality predicts outcomes
\end{itemize}

\subsection{Synthesis: Heckman's Contribution to the Framework}

\paragraph{The Central Insight.}
Heckman's work establishes that human capital formation exhibits
\emph{dynamic complementarity}: skills beget skills, and early investments
multiply the returns to later investments.

\paragraph{The Heckman Equation.}
\begin{equation}
  \frac{\partial^2 \theta_{t+1}}{\partial \theta_t \partial I_t} > 0
\end{equation}
This positive cross-partial is the foundation for understanding why
early childhood investments have such high returns.

\paragraph{Five Major Contributions.}

\begin{enumerate}
  \item \textbf{Selection correction:} Heckman correction for sample selection
  \item \textbf{Treatment effects:} Heterogeneous responses to policy
  \item \textbf{Skill formation:} Technology of human development
  \item \textbf{Multiple skills:} Cognitive + noncognitive matter
  \item \textbf{Early investment:} High returns, no equity-efficiency tradeoff
\end{enumerate}

\paragraph{Framework Elements Established.}

\begin{center}
\renewcommand{\arraystretch}{1.2}
\small
\begin{tabular}{lll}
\hline
\textbf{Element} & \textbf{Papers} & \textbf{Contribution} \\
\hline
Selection bias & 1, 2, 3 & Context determines observation \\
Treatment effects & 5, 6, 14 & Heterogeneous policy responses \\
Dynamic complementarity & 7, 8, 9 & Skills beget skills \\
Multiple skills & 10, 11, 20 & Cognitive + noncognitive \\
Early investment & 12, 13 & High ROI, equity + efficiency \\
General equilibrium & 16, 17 & Policy changes $\Psi$ \\
\hline
\end{tabular}
\end{center}

\paragraph{Connection to Other Research Programs.}

\begin{center}
\renewcommand{\arraystretch}{1.2}
\begin{tabular}{lllll}
\hline
& \textbf{Fehr} & \textbf{Acemoglu} & \textbf{Shleifer} & \textbf{Heckman} \\
\hline
Focus & Preferences & Political inst. & Legal inst. & Human capital \\
$C$ insight & Social pref. & Network & Cognitive & Dynamic compl. \\
$\Psi$ type & Norms & Institutions & Legal origins & Family, school \\
Time scale & Sessions & Centuries & Decades & Life cycle \\
Nobel & (behavioral) & 2024 & (with LLSV) & 2000 \\
\hline
\end{tabular}
\end{center}

\paragraph{The Four Pillars United.}
The $(C, \Psi)$ framework unifies:
\begin{itemize}
  \item \textbf{Fehr:} Preferences shaped by social context ($\Psi_{norm}$)
  \item \textbf{Acemoglu:} Institutions as fundamental cause ($\Psi_{inst}$)
  \item \textbf{Shleifer:} Legal origins determine markets ($\Psi_{legal}$)
  \item \textbf{Heckman:} Skills formed through complementary investments ($C^{skill,investment}$)
\end{itemize}

Together, these four research programs establish that context ($\Psi$)
determines equilibrium ($C^*$), and complementarity structure ($C$)
determines how interventions propagate through the system.

%%%%%%%%%%%%%%%%%%%%%%%%%%%%%%%%%%%%%%%%%%%%%%%%%%%%%%%%%%%%%%%%%%%%%%%%%%%%%%%


%%%%%%%%%%%%%%%%%%%%%%%%%%%%%%%%%%%%%%%%%%%%%%%%%%%%%%%%%%%%%%%%%%%%%%%%%%%%%%%
%%%%%%%%%%%%%%%%%%%%%%%%%%%%%%%%%%%%%%%%%%%%%%%%%%%%%%%%%%%%%%%%%%%%%%%%%%%%%%%

%%%%%%%%%%%%%%%%%%%%%%%%%%%%%%%%%%%%%%%%%%%%%%%%%%%%%%%%%%%%%%%%%%%%%%%%%%%%%%%
\section{Key Results}
\label{sec:hec-results}
%%%%%%%%%%%%%%%%%%%%%%%%%%%%%%%%%%%%%%%%%%%%%%%%%%%%%%%%%%%%%%%%%%%%%%%%%%%%%%%

The analysis yields the following key results:

\begin{enumerate}
    \item \textbf{Result 1:} [Primary finding with quantitative support]
    \item \textbf{Result 2:} [Secondary finding with implications]
    \item \textbf{Result 3:} [Practical application or boundary condition]
\end{enumerate}


\section{Summary}
\label{sec:hec-summary}
%%%%%%%%%%%%%%%%%%%%%%%%%%%%%%%%%%%%%%%%%%%%%%%%%%%%%%%%%%%%%%%%%%%%%%%%%%%%%%%

This appendix has presented the key concepts and theoretical foundations.
The main contributions include:

\begin{enumerate}
    \item Formal definitions and theoretical framework
    \item Integration with the broader EBF structure
    \item Practical guidelines for application
\end{enumerate}

The content supports decision-making and behavioral analysis within the
Evidence-Based Framework, providing a foundation for further research
and practical implementation.



%%%%%%%%%%%%%%%%%%%%%%%%%%%%%%%%%%%%%%%%%%%%%%%%%%%%%%%%%%%%%%%%%%%%%%%%%%%%%%%
\section{Critical Foundations and Objections}
\label{sec:hec-foundations}
%%%%%%%%%%%%%%%%%%%%%%%%%%%%%%%%%%%%%%%%%%%%%%%%%%%%%%%%%%%%%%%%%%%%%%%%%%%%%%%

\subsection{Objection 1: Scope Limitations}

\textbf{Objection:} The framework presented may have limited applicability
outside the specific domain contexts discussed.

\textbf{Response:} We acknowledge domain-specificity. The framework provides
general principles that should be calibrated to specific contexts. Cross-domain
validation remains an open research question.

\subsection{Objection 2: Measurement Challenges}

\textbf{Objection:} Key constructs may be difficult to measure reliably in
practice.

\textbf{Response:} We recommend using validated scales and instruments where
available, and developing domain-specific proxies where necessary. Sensitivity
analysis should assess robustness to measurement uncertainty.



%%%%%%%%%%%%%%%%%%%%%%%%%%%%%%%%%%%%%%%%%%%%%%%%%%%%%%%%%%%%%%%%%%%%%%%%%%%%%%%
\section{Glossary of Symbols}
\label{sec:hec-glossary}
%%%%%%%%%%%%%%%%%%%%%%%%%%%%%%%%%%%%%%%%%%%%%%%%%%%%%%%%%%%%%%%%%%%%%%%%%%%%%%%

For comprehensive definitions, see Appendix G (Master Glossary).

\begin{table}[htbp]
\centering
\caption{Key Symbols in Appendix UNMAPPED_HEC}
\begin{tabular}{lll}
\toprule
\textbf{Symbol} & \textbf{Meaning} & \textbf{Reference} \\
\midrule
$\gamma$ & Complementarity parameter & Appendix UNMAPPED_B \\
$\Psi$ & Context vector & Appendix UNMAPPED_V \\
$U$ & Utility function & Chapter 3 \\
\bottomrule
\end{tabular}
\end{table}


\section{References}
\label{sec:hec-references}
%%%%%%%%%%%%%%%%%%%%%%%%%%%%%%%%%%%%%%%%%%%%%%%%%%%%%%%%%%%%%%%%%%%%%%%%%%%%%%%

See master bibliography (\texttt{bcm\_master.bib}) for complete citations.

\nocite{bcm_master}

